% META: {"id": "1NSI-PM-CO-021", "chapitre": "1NSI-PROJET-METHODES", "type_objet": "correction", "exercice_ref": "1NSI-PM-EX-003", "capacites": ["1NSI-PROJET-METHODES-C4"], "sources_inspiration": [], "status": "approved", "capacites_codes": ["C4"], "parcours": 1, "fichier_exercice": "chapitres/1NSI-PROJET-METHODES/exercices/1NSI-PM-EX-021.tex"}
% BEGIN-VERIFY
% notes = {"Alice": 15, "Bob": 12, "Chloe": 18}
% def afficher_note_corrigee(notes, nom):
%     if nom in notes:
%         return f"{nom} : {notes[nom]}"
%     else:
%         return "Nom inconnu"
% assert afficher_note_corrigee(notes, "David") == "Nom inconnu"
% assert afficher_note_corrigee(notes, "Alice") == "Alice : 15"
% END-VERIFY
\begin{corrige}{1NSI-PM-EX-003}
\textbf{1.} Le programme lève une \lstinline{KeyError}, avec le message
\lstinline{'David'}.

\textbf{2.} La ligne fautive est \lstinline{notes[nom]} dans \lstinline{afficher_note} :
elle tente d'accéder à la clé \lstinline{"David"} dans le dictionnaire \lstinline{notes},
mais cette clé n'existe pas (seuls \lstinline{"Alice"}, \lstinline{"Bob"} et
\lstinline{"Chloe"} y figurent). La cause est donc un appel de la fonction avec un nom
absent du dictionnaire.

\textbf{3.}
\begin{python}
def afficher_note_corrigee(notes, nom):
    if nom in notes:
        return f"{nom} : {notes[nom]}"
    else:
        return "Nom inconnu"
\end{python}
\end{corrige}
